\chapter{Literature Review}
\label{ch:2}
\label{sec:2}


\section{Comparison Framework}
\label{sec:2_1}


For the large-scale application of archetype-based UBEM, previous studies can be organised into four workflow routes according to how building information is transferred from observational sources to energy-related tasks: the reference-attribute, proxy-attribute, direct visual, and attribute-mediated visual routes. These routes rely on different forms of data availability and differ in whether they retain explicit building attributes and the TABULA archetype structure. \cref{tab:2_1} establishes a common comparison framework based on the observed input, intermediate representation, typical target, and principal dependency of each route.


\begin{landscape}
\begin{table}[p]
\centering
\caption{Comparison framework for literature-review routes}
\label{tab:2_1}
\small
\begin{tblr}{
  width = \linewidth,
  colspec = {X[1.7,l] X[1,l] X[1,l] X[1,l] X[1,l]},
  row{1} = {font=\bfseries},
  rowsep = 3pt,
  colsep = 4pt
}
\toprule
Route & Observed input & Intermediate representation & Typical target & Principal dependency \\
\midrule
Reference-attribute & Registry, geometry, and neighbourhood data & Recorded building attributes & EPC or energy-performance prediction & Complete and linkable building-level records \\
Proxy-attribute & Aggregated census, morphology, or remote-sensing data & Inferred attributes or archetype classes & Archetype assignment or energy modelling & Transferable relationship between the proxy and the individual building \\
Direct visual & SVI, aerial, or multimodal imagery & Latent visual representation & Direct EPC classification & Local labels, image coverage, and domain stability \\
Attribute-mediated visual & SVI & Explicit building type, construction year, floor count, and TABULA cell & Archetype parameters and downstream EPC classification & Attribute extractability and error propagation \\
\bottomrule
\end{tblr}
\end{table}
\end{landscape}


The studies considered here differ in their label spaces, source populations, data splits, and metrics. Their numerical results are therefore used to characterise the capabilities and boundaries of each route, rather than to rank performance directly across studies.


\section{Reference and Open-Data Approaches}
\label{sec:2_2}


When complete building-level records are available, the reference-attribute route uses administrative, geometric, and neighbourhood data to train statistical models. \hyperlink{ref:hettinga2023}{Hettinga et al. (2023)} developed an EPC classifier for Dutch residential buildings by combining BAG building year, building type, ground-floor area, total floor area, and exposed-envelope geometry with CBS neighbourhood-level home ownership and average building value. Their random forest achieved 71\% accuracy and 93\% \ensuremath{\pm}1-grade accuracy on the held-out test subset of an 80/10/10 train--validation--test protocol. By comparison, TABULA rule-based assignment achieved 26\% accuracy and 67\% \ensuremath{\pm}1-grade accuracy across 2,574,227 EPC records. The two approaches used different evaluation sets, and the EPC reference labels covered A++--G whereas the TABULA outputs covered only B--G. These values therefore demonstrate the analytical potential of additional open-data features, rather than providing a fully equivalent model comparison.


This route retains interpretable tabular inputs and can incorporate building-level and neighbourhood-level information beyond the archetype. However, it assumes that the target region already has linkable reference attributes of sufficient quality. When building type, construction year, or floor count is missing, the route alone cannot determine how these inputs should be reconstructed or quantify the cost of replacing them with information from other sources.


\section{Proxy-Based Attribute and Archetype Inference}
\label{sec:2_3}


The proxy-attribute route uses data sources that correlate with a target attribute but do not constitute observations of the individual building itself. \hyperlink{ref:dabrock2025}{Dabrock et al. (2025)} combined German census data aggregated to a 100 m \ensuremath{\times} 100 m grid with building-level and neighbourhood-level morphological and socioeconomic features, and trained XGBoost models to predict residential size class and construction-year class. The two tasks achieved accuracies of 97.4\% and 73.9\%, respectively, on an individual-building test set, and the predictions were subsequently used to construct a nationwide TABULA-based building dataset. The study also observed that the resulting dataset tended to overrepresent majority classes, indicating that aggregate proxies can support individual-building estimates but that their errors are affected by grid composition and class imbalance.


Remote sensing can also serve as an attribute proxy. \hyperlink{ref:wurm2021}{Wurm et al. (2021)} extracted construction type and construction period from aerial imagery and linked these attributes to urban heat-demand modelling. Their study demonstrated that imagery-derived attributes can be incorporated into a downstream energy model. However, the imagery source was not SVI, and reference attributes and image-derived substitutes were not compared for the same buildings under a common downstream protocol. Proxy-based methods therefore offer a feasible route for large-scale data completion, but their transferability remains dependent on proxy resolution, geographic context, and target taxonomy.


\section{Direct Image-to-EPC Approaches}
\label{sec:2_4}


The direct visual route bypasses explicit archetype assignment and predicts EPC labels directly from imagery. \hyperlink{ref:mayer2023uk}{Mayer et al. (2023)} integrated SVI, aerial imagery, building footprints, and land-surface temperature for approximately 40,000 buildings in the United Kingdom and performed binary classification of A--D versus E--G. The complete multimodal model achieved a macro-F1 score of 64.64\%. In a study of more than 21,000 buildings in the Copenhagen metropolitan area, the same research group obtained a macro-F1 score of 58.79\% using SVI and aerial imagery, and concluded that these two image sources alone were insufficient to distinguish building energy efficiency consistently \hyperlink{ref:mayer2023dk}{(Mayer, Heilborn, \& Fischer, 2023)}.


Subsequent studies expanded both the data sources and label granularity. \hyperlink{ref:liu2025}{Liu et al. (2025)} combined Google Street View and GIS data in a three-stage framework using 1\% local samples for transfer learning. They reported an accuracy of 84\% for A--G EPC classification, although performance decreased for occluded high-rise buildings. \hyperlink{ref:sun2026}{Sun et al. (2026)} integrated thermal infrared imagery, optical remote sensing, SVI, morphology, and socioeconomic indicators. Their binary high- versus low-performance tasks achieved F1 scores of 0.64 in Glasgow and 0.69 in Edinburgh. In their ablation study, the SVI-only model achieved F1 scores of 0.59 and 0.68 in the two cities, respectively.


These findings support the presence of EPC-predictive signals in building appearance and urban-context data, but the capability of the direct visual route cannot be summarised by a single performance value. Previous studies differ substantially in modality, local adaptation, label aggregation, class definition, and evaluation protocol. Most employ binary targets, while only a small number report A--G classification; consequently, a directly comparable full-grade benchmark has not yet emerged. Moreover, direct models bypass explicit TABULA parameters. Although they can directly support classification, they cannot independently identify the contribution of individual building attributes or directly quantify how attribute substitution affects the archetype cell and its inherited physical parameters.


\section{Image-Based Building-Attribute Extraction}
\label{sec:2_5}


The attribute-mediated visual route first infers explicit building attributes from SVI and then performs TABULA cell assignment or downstream modelling. In contrast to the direct route, this structure retains intermediate attributes and thereby enables separate evaluation of attribute extraction, joint-cell assignment, and downstream utility. The relevant vision strategies can be further divided into supervised full fine-tuning, frozen self-supervised representations, and inference without parameter updating.


\subsection{Supervised Full Fine-Tuning}
\label{sec:2_5_1}


OpenFACADES compared CNN, Vision Transformer, and VLM configurations using 31,180 labelled images from seven cities \hyperlink{ref:liang2025}{(Liang et al., 2025)}. Fine-tuned InternVL3-2B achieved an accuracy of 61.27\% for building-type classification. For building-age regression, it achieved an R\textsuperscript{2} of 0.710 and an MAE of 28.06 years, while the corresponding values for floor-count regression were 0.789 and 2.13 floors. The study also applied the model to large-scale building-profile enrichment, demonstrating that a single image-processing framework can support multi-attribute inference. Full fine-tuning, however, still requires a labelled training corpus, and its applicability across cities or taxonomies must be established through independent transfer evaluation.


\subsection{Frozen Self-Supervised Representations}
\label{sec:2_5_2}


Self-supervised visual foundation models learn general-purpose representations from large-scale image corpora without manual annotation. DINOv2 provides visual features that can be used with task-specific prediction heads without updating the backbone parameters \hyperlink{ref:oquab2023}{(Oquab et al., 2023)}. Compared with full fine-tuning, a frozen-backbone configuration reduces the number of parameters that must be updated and separates representation learning from downstream supervised fitting. However, the building-energy studies reviewed above do not provide controlled evidence that compares a frozen self-supervised backbone with a fully fine-tuned CNN and VLM inference on the same building population and traces their attribute errors through TABULA cell assignment to EPC classification.


\subsection{VLM Inference Without Parameter Updating}
\label{sec:2_5_3}


VLMs offer another route for building-attribute inference without updating model parameters. \hyperlink{ref:zeng2024}{Zeng et al. (2024)} used GPT-4 Vision for zero-shot classification of building-age epochs from London façade images, achieving an accuracy of 39.69\% and an epoch-class MAE of 0.85. The model was less reliable for very old buildings and for fine-grained distinctions between periods separated by no more than two decades. OpenFACADES also compared zero-shot and fine-tuned VLM configurations, with its overall results showing that fine-tuning improved multi-attribute performance \hyperlink{ref:liang2025}{(Liang et al., 2025)}. Inference without parameter updating therefore reduces the need for task-specific training, but its results remain sensitive to model access, prompt design, output constraints, and geographic context. Furthermore, the term ``zero-shot'' does not necessarily mean that development examples played no role in prompt selection.


\section{Archetype-Grounded and Semi-Synthetic Datasets}
\label{sec:2_6}


The routes reviewed above also produce different forms of building-level datasets. \hyperlink{ref:dabrock2025}{Dabrock et al. (2025)} mapped predicted size and construction-year classes to TABULA archetypes to construct a nationwide residential dataset. OpenFACADES, by contrast, extended image-derived attributes and captions to approximately 1.2 million images covering 500,000 buildings \hyperlink{ref:liang2025}{(Liang et al., 2025)}. These studies demonstrate archetype instantiation and image-based data enrichment, respectively, but neither organises reference attributes, image-derived substitutes, deterministically assigned TABULA parameters, and EPC reference labels within a single paired-comparison structure.


The recent Building Energy Graphs (BEGs) preprint further combined OpenStreetMap footprints, GHSL products, European building archetypes, and deterministic assumptions to construct simulation-ready building energy models, with EnergyPlus outputs supporting GNN surrogate prediction \hyperlink{ref:dalach2025}{(Dalach et al., 2025)}. This work demonstrates that observed geospatial sources, archetype assumptions, and simulation outputs can be integrated into synthetic or semi-synthetic learning artefacts. Its target, however, is simulated energy-use intensity rather than administrative EPC reference labels. This distinction positions the utility of the dataset developed in the present study as measurement-anchored attribute-substitution and EPC analysis, rather than as another simulation-surrogate benchmark.
